LLM agent pipelines routinely delegate \emph{per-item judgments} to a model in batches: given a policy document
and a list of records, judge every record in one call. We report a failure mode of this pattern discovered in a
live tool-use pipeline ($\tau^2$-bench banking): specific items are judged incorrectly \emph{only when batched
with other items that engage the same policy clause}, and the failure is a silent retreat to a default answer,
not a random error. Through a chain of controlled probes on the live prompt (temperature~0, free-form
generation) we exclude token load, absolute context position, generation order, and category similarity, and
isolate the interference source: \textbf{previous items that inference-bind the same rule clause without any
lexical name match}. Items anchored by explicit lexical evidence are immune; judgment-dependent items fail with
a sharp threshold at $k^\ast{=}2$ interfering items, with output instability (unit slips such as
$5 \rightarrow 500$) at the threshold. A double dissociation rejects quota-consumption, judgment-budget, and
category-cue accounts. The mechanism-derived mitigation---capping judgment batches at 2 items so every item
sits at a list edge---closes the failure in probes (38/38), causes no regression across all probed cells
(166/167), and turns previously failing live tasks into passes, while prompt mitigations fail entirely,
consistent with prior reports. A logit-readout size sweep (1.5B--32B) shows the staircase is not an artifact of
one model scale, and a structural replication in a second domain (airline baggage fees, prompts under 1.5k
tokens) shows it is not an artifact of one domain or of long context. We position the phenomenon against
serial-position and proactive-interference literature: what is new is the \emph{unit} of interference (the rule
clause), the dissociation methodology, the default-retreat failure mode, and a structural fix derived from the
mechanism rather than from folklore.
